\section{The construction used as input}\label{sec:construction}

We record precisely the part of Stolarski's construction that is needed.
The role of this input is to produce a single smooth singular flow with
compatible inner signs and cone profiles. The curvature estimates proved
in the later sections will be consequences of these data.
Throughout, set \(p=2\) and write
\begin{equation}\label{eq:parameters}
\begin{gathered}
m=p+q,\qquad A^2=\frac{p-1}{m-1},\qquad B^2=\frac{q-1}{m-1},\\
d=\frac{m-1-\sqrt{(m-1)(m-9)}}2,\qquad
\nu=k-d/2,\qquad \sigma=k/d.
\end{gathered}
\end{equation}
Here \(q\ge10\), and \(k\) is a sufficiently large even positive integer.
In particular \(\nu>0\). We take \(A,B>0\) and use the time-dependent scales
\begin{equation}\label{eq:scales}
\delta=T-t,\qquad \tau=-\log\delta,\qquad
\theta=\delta^\sigma,\qquad \gamma=r/\sqrt\delta.
\end{equation}
The dimension of the manifold is \(m+1\).
The identities \(\nu/d=\sigma-1/2\) and
\(\sqrt\delta\,\delta^{\nu/d}=\theta\) relate the two scales.

\begin{proposition}[Construction input]\label{prop:input}
For the parameters above, Stolarski's construction supplies a smooth Ricci
flow up to a finite time \(T\) on \(\Sp^p\times\Sp^{q+1}\) with the
following properties, after translating its starting time to zero.

\begin{enumerate}[label=\textup{(\roman*)}]
\item Off the two pole orbits it has the form
\[
g=ds^2+\phi^2g_{\Sp^p}+r^2g_{\Sp^q}.
\]
Here \(s\) is radial arclength at each time. The metric is invariant under
reflection exchanging the poles. On each open hemisphere, oriented from
its pole toward the equator, \(r_s>0\). At a pole, \(r=0\), \(r_s=1\),
and \(\phi_s=0\); the usual odd and even polar expansions hold for \(r\)
and \(\phi\), respectively, and \(\phi>0\) for every \(t<T\).

\item Write \(u=\log\phi\), \(z=r_s^2\), and
\[
\Ut=\log\frac{\phi}{(A/B)r},\qquad \Zt=z-B^2.
\]
There are positive tolerances \(\eta_j^U,\eta_j^Z\), large fixed numbers
\(1<Y_1\le Y_2\), and constants \(M_0>0\), \(0<\beta<1/2\), such that
for \(j=0,1,2\),
\begin{equation}\label{eq:profile-bounds}
\begin{aligned}
\left|\partial_\gamma^j(\Ut-\delta^\nu U_k)\right|
&\le\eta_j^U\delta^\nu
       \bigl(\gamma^{-d-j}+\gamma^{2k-d}\bigr),\\
\left|\partial_\gamma^j(\Zt-\delta^\nu Z_k)\right|
&\le\eta_j^Z\delta^\nu
       \bigl(\gamma^{-d-j}+\gamma^{2k-d}\bigr).
\end{aligned}
\end{equation}
The first estimate holds from \(Y_1\delta^{\nu/d}\) to
\(M_0e^{\beta\tau}\), and the second from \(Y_2\delta^{\nu/d}\) to that
upper endpoint, within the hemisphere. The profiles are smooth on
\(\gamma>0\). The function \(U_k\) is \(\gamma^{-d}\) times a polynomial
of degree \(k\) in \(\gamma^2\), with positive leading coefficients at
zero and infinity:
\[
U_k(\gamma)\sim c_1\gamma^{-d}\quad(\gamma\downarrow0),\qquad
U_k(\gamma)\sim c_2\gamma^{2k-d}\quad(\gamma\to\infty),\qquad c_1,c_2>0.
\]
Also \(Z_k(\gamma)=O_{q,k}(\gamma^{-d})\) as \(\gamma\downarrow0\).

\item For all sufficiently late times, putting \(f=ru_r\), one has
\begin{equation}\label{eq:inner-inputs}
\begin{gathered}
\phi/r\ge A/B,\qquad 0\le f\le1,\qquad f_r\ge0
       \quad(0<r\le Y_1\theta),\\
B^2\le z\le1\quad(0<r\le Y_2\theta).
\end{gathered}
\end{equation}
These inner intervals lie strictly below the equator.

\item At the initial time \(\tau_0=-\log T\), the same weighted profile
bounds hold up to an initial outer cutoff
\(G_{\mathrm{init}}=M_0e^{\beta_{\mathrm{init}}\tau_0}\),
where \(0<\beta_{\mathrm{init}}<1/2\), and \(\Ut\ge0\) beyond that
cutoff. The exponents are chosen in the ranges permitted by the
construction; in particular, one may take
\[
0<\beta<\frac{\nu}{2\nu+1}\le\beta_{\mathrm{init}}<\frac12.
\]
They are therefore distinct choices, not two arbitrary exponents.
The tolerances can be
chosen sufficiently small depending on \(q,k\). The inner cutoffs can then
be increased in order, first \(Y_1\) and then \(Y_2\); the initial
rescaled time can subsequently be increased. Each such increase is
subject only to lower-size requirements from the construction.
\end{enumerate}
\end{proposition}

\begin{proof}[Source and change of notation]
The existence statement is the solution furnished by Lemma~3.12 and the
proof of Theorem~1.1 in~\cite{Stolarski2019}. The weighted estimates are
Definition~3.1. The initial family is Definition~3.5 and Lemma~3.6,
with the outer conditions of Definition~3.2.
The inner inequalities are the Inner Region Barriers I in Definition~3.3,
as propagated by Lemmas~4.5--4.9 and explicitly invoked in the proof of
Theorem~4.3 on page~25. In particular that proof gives the barrier
conditions on the inner intervals; the shorter list in Definition~4.1
of its final class \(\mathcal P\) is not the only information being used.
The eigenprofile properties are Propositions~A.12, A.13, and A.21;
the parameter order is listed in Appendix~B.

For completeness, the initial weighted bounds hold up to the initial
expansion's own upper endpoint, not just the smaller propagated endpoint.
The lower \(U\)-modes in Definition~3.5 are \(\gamma^{-d}\) times
polynomials in \(\gamma^2\) of degree below \(k\).
By Propositions~A.10 and~A.19, the homogeneous lower \(Z\)-modes are
\(\gamma^2\) times polynomials, with growth exponents strictly below
\(2k-d\). Each of these finitely many modes and its first two
derivatives is therefore bounded by the corresponding weight in
\eqref{eq:profile-bounds} on all \(\gamma>0\).
Taking the lower-mode coefficient bound at the initial time to be
\(\epsilon_0\delta^\nu\), with \(\epsilon_0\) sufficiently small,
gives the required tolerances. This choice is compatible with the degree
argument in Lemma~3.10: the projection errors in Lemmas~6.13--6.14 can
subsequently be reduced by increasing the initial rescaled time.
Definition~3.5 imposes the outer conditions starting at that same initial
expansion endpoint. Thus the initial profile bounds and outer sign meet
at \(G_{\mathrm{init}}\); no identification of
\(\beta_{\mathrm{init}}\) with \(\beta\) is required.
The relation between their permitted ranges is specified in the proof
of Theorem~4.3, page~26, of~\cite{Stolarski2019}.

In the source's notation,
\(B^2\lambda_k=-\nu\), \(\alpha_k=\nu/d\), and
\(Y_1=\Upsilon_U\), \(Y_2=\Upsilon_Z\).
Thus \(\sqrt\delta\,Y_i e^{-\alpha_k\tau}=Y_i\theta\).
Moreover
\[
f=1+\gamma\Ut_\gamma,\qquad
f_r=\frac{\gamma}{\sqrt\delta}
   \left(\Ut_{\gamma\gamma}+\frac{\Ut_\gamma}{\gamma}\right),
\]
which translates the source's log-radius convexity into the last slope
inequality in~\eqref{eq:inner-inputs}.
The mode index \(k\) here is the index in the construction, rather than
the arbitrary curvature blow-up exponent in the statement of its
Theorem~1.1.
\end{proof}

\begin{remark}\label{rem:input-scope}
Proposition~\ref{prop:input} is the external existence input to this paper.
We do not import a scalar-curvature upper bound, a full-curvature bound at
the cap scale, or convergence to a particular complete cap metric.
The estimates establishing those curvature bounds are given below.
\end{remark}
